% Transcription — fonds Alexandre Grothendieck, shelfmark 57, batch 12 (pages 221-240).
% Established from the facsimile archives/batches/57/batch-12.pdf, per the skill
% transcribe-grothendieck.
%
% Pass: Opus 5.5 (claude-opus-5-5), 2026-10-03 — first pass, unchecked against the pages by a human.
%
% Pages transcribed: 221, 223, 226-230, 232, 234, 236, 238, 240; page 225 (a cover
% in his hand) gives the section heading and has no \page.
% Absent: 222, 224 (typescript versos of a Bourbaki draft, n° 215, upside down),
% 231 (blank), 233, 235, 237, 239 (typescript versos of a Bourbaki draft
% « Algèbre commutative », n° 227 and n° 338).
%   221-223: the family xy(x+y+tz) = 0 over Z[t]: flatness, the three components,
%     H^1(X, O_X) by descent along X_1 + X_2 + X_3 -> X.
%   226-230: his typed letter to Hironaka, Neuilly, 6 July 1962, with his ink
%     corrections: Mittag-Leffler for Pic(X_n), formal and analytic local
%     Picard groups, conjectures for henselian "good" rings.
%   232-240: Pic of a formal scheme in char. 0: a lemma on pi_0 of homomorphisms
%     of group schemes locally of finite type, corollaries on projective systems
%     (G_n) and Mittag-Leffler, application to (Pic_{X_n/k}).
\documentclass[11pt,a4paper]{article}
\input{../preamble/grothendieck}

\folder{57}
\batch{12}
\pages{221}{240}
\foldertitle{Picard : tapuscrit annoté (s.d.), lettres (s.d., 1962).}
\dating{1962-[vers 1968]}
\watermark{Édition de démonstration}

\begin{document}

\page{221}
\note{la page s'ouvre sur un exemple déjà en cours ; les notations $X$, $S$, $X^{(n)}$, $S^{(n)}$ viennent de pages antérieures au lot.}

\begin{tikzcd}
X^{(n)}_{a} \arrow[d] & \mathcal{S} \arrow[d] \\
S^{(0)}_{e} = \operatorname{Spec}(\mathbf{Z}[t]) & S \arrow[l]
\end{tikzcd}

$t \in \Gamma(S, \underline{O}_S)$,
\quad \struck{$A[x,y]/xy(x+y+t)$}
\[
A[x,y,z]/xy(x+y+tz) = \mathcal{S}
\]
\note{au-dessus de la flèche verticale de gauche, un symbole entièrement noirci.}

\struck{$xy(x+y+tz)$}
\[
\begin{array}{ccc}
(x) & (y) & (x+y+tz) \\
\mathfrak{p}_1 & \mathfrak{p}_2 & \mathfrak{p}_3 \\
\underline{I}_1 & \underline{I}_2 & \underline{I}_3 \\
X_1 & X_2 & X_3
\end{array}
\]

\textbf{Lemme 1} $X$ est plat sur $S$.

Il suffit de prouver que $X^{(n)}$ est plat sur $S^{(n)}$, donc
\uncertain{à fortiori} \struck{\ill{}} $k[x,y,z]/xy(x+y+tz)$ est plat sur $k[t]$,
donc sans torsion. On a \uncertain{supposé} que $\varphi \in k[t]$,
\add{$\varphi \neq 0$}, $F \in k[x,y,z,t]$ tels que
$\varphi F \in xy(x+y+tz)$, donc $\varphi F \in (x)$, $\varphi F \in (y)$,
$\varphi F \in (x+y+tz)$, d'où comme \struck{$k[x,y,z,t]$} les idéaux
\struck{$(x)$, $(y)$, $(x+y+tz)$}
$\mathfrak{p}_1^{(0)} = (x)$, $\mathfrak{p}_2^{(0)} = (y)$,
$\mathfrak{p}_3^{(0)} = (x+y+tz)$ sont

\marginal{\struck{Lemme \ill{}}} premiers, et
$k[t] \to k[x,y,z,t]/\mathfrak{p}_i^{(0)}$ est injectif,

\marginal{\struck{Lemme \ill{}}} la relation $F \in$ \struck{$(x)$, \ill{}}
$\mathfrak{p}_i^{(0)}$ pour \uncertain{$i = 1,2,3$}, d'où
$F \in \mathfrak{p} = xy(x+y+tz)$ \ill{} \ill{}.
\note{dans la marge, deux étiquettes « Lemme » surchargées et biffées, dont le numéro ne se lit pas : la numérotation des lemmes est en train de changer.}

\textbf{Lemme 0} On a $\mathfrak{p}_1 \cap \mathfrak{p}_2 \cap \mathfrak{p}_3 = 0$
i.e. $(x) \cap (y) \cap (x+y+zt) = (xy(x+y+zt))$.
\note{la fin de la ligne est coupée par le bord du feuillet ; le membre de droite est lu d'après la page suivante.}

\page{223}
(\emph{N.B.} Ce dernier lemme est vrai pour $X/S$.)
Soit en effet $F \in (x) \cap (y) \cap (x+y+zt)$, on a donc $F = xyG$, et
$F$ \uncertain{devient nul} \uncertain{aussi} \ill{} $y$ \uncertain{faisant}
$y = -(x+zt)$, d'où $x(x+zt)\,G(x,-x-zt,z) = 0$, or $x$ et $x+zt$ sont
non diviseurs de $0$ dans
$\underline{O}_S[x,y,z] = \underline{O}_S[y,z][x]$ \struck{\ill{}}, d'où
$G(x,-x-zt,z) = 0$ i.e.\ $G$ \add{devient} \struck{\ill{}} nul pour la
substitution $y = -x-zt$, donc $G$ est divisible par $y+x+zt$.

\emph{Corollaire du lemme 0} $\underline{I}_1 \cap \underline{I}_2 \cap \underline{I}_3 = (0)$, donc
\[
\bigvee_{1 \leqslant i \leqslant 3} X_i = X .
\]

\textbf{Lemme 2} \struck{Si $i \neq j$, on a $\mathcal{S}/\mathfrak{p}_i + \mathfrak{p}_j \simeq \underline{O}_S[u]$}

Pour $i = 1,2,3$, $\mathcal{S}/\mathfrak{p}_i \simeq \underline{O}_S[u,v]$
donc $X_i \simeq \mathbf{P}^1_S$ ;

Si $i \neq j$, $\mathcal{S}/\mathfrak{p}_i + \mathfrak{p}_j \simeq \underline{O}_S[u]$
\struck{et} donc $X_{ij} \simeq S$ ;

\struck{Si} \add{Enfin} $\mathcal{S}/\mathfrak{p}_1 + \mathfrak{p}_2 + \mathfrak{p}_3 \simeq (\underline{O}_S/t\underline{O}_S)[u]$
donc $X_{123} \simeq (S, \underline{O}_S/t) = t^{-1}(0)$.
\note{les trois énoncés sont réunis par une accolade ; « $\mathbf{P}^1_S$ » et « $X_{ij} \simeq S$ » sont lus tels qu'écrits.}

\emph{Corollaire} Si $S$ est affine, on a
\[
H^1(X, \underline{O}_X) \simeq H^1(X_1 \sqcup X_2 \sqcup X_3 / X, \underline{G}_a) \simeq
\]
\marginal{justification : \struck{\ill{}} $X_1 \sqcup X_2 \sqcup X_3 \to X$ est un \uncertain{morphisme} de \uncertain{descente} \uncertain{stricte} pour les \uncertain{schémas}}
\note{la phrase s'arrête sur « $\simeq$ » en bas de page ; la suite n'est pas dans le lot.}

\section{Schémas de Picard des schémas formels en car.~0, ou des schémas propres sur une algèbre analytique sur~$\mathbf{C}$}
\note{titre écrit de sa main seul sur la page 225, qui sert de chemise à la lettre qui suit.}

\page{226}
\note{lettre de Grothendieck à H.~Hironaka, dactylographiée en anglais, avec ses corrections à l'encre. La machine n'a pas frappé certains symboles ($\to$, $\in$, $\ni$) ; ils sont restitués entre crochets d'éditeur quand le texte les fixe.}

Neuilly July 6 1962

Dear Hironaka,

I had a little thought over our conversation last tuesday, it occured to me
that the type of argument I used yields in fact the following stronger result:

\emph{Theorem} Let $f\colon X$ \supplied{$\to$} $Y$ be a proper morphism of
analytic spaces over $\underline{C}$, let $y$ \supplied{$\in$} $Y$,
$Y_n = \operatorname{Spec\,an}(\underline{O}_y/\underline{m}_y^{n+1})$,
$X_n = X \times_Y Y_n$,
\begin{align*}
\operatorname{Pic}(X_y) &= R^1 f(\underline{O}_X^{\bullet})_y
  = \varinjlim_{U \ni y} \operatorname{Pic}(f^{-1}(U)) ,\\
\operatorname{Pic}(\hat{X}_y) &= \varprojlim \operatorname{Pic}(X_n) ,
\end{align*}
and consider the canonical homomorphisms
\[
\operatorname{Pic}(X_y) \xrightarrow{\ u\ } \operatorname{Pic}(\hat{X}_y)
\xrightarrow{\ v_n\ } \operatorname{Pic}(X_n)
\]
Then the following are true:

(i) The inverse system $(\operatorname{Pic}(X_n))_{n \in \underline{N}}$
satisfies the condition of Mittag-Leffler (even with Artin-Riesz type of
uniformity).

(ii) $\operatorname{Im} v_n = \operatorname{Im} v_n u$ $=$ (in virtue of (i))
set of universal images of $\operatorname{Pic}(X_n)$ in the inverse system
$(\operatorname{Pic}(X_m))_{m \in \underline{N}}$.

(iii) In order for $u$ to be an isomorphism, it is nec and suff that
$R^1 f_*(\underline{O}_X)$ has a support not containing $y$ or having $y$ as
an isolated point, i.e.\ $R^1 f_*(\underline{O}_X)_y$ is a module of finite
length.

In fact (i) can be made more precise:

(i bis) In the inverse system
$(\underline{\mathrm{Pic}}_{X_n/\underline{C}})_{n \in \underline{N}}$ of
analytic groups, the system of the “Néron-Sévéri groups” is constant for $n$
large, whereas for $m \geqslant n$ and $n$ large, the Kernel and Cokernel of
\[
\underline{\mathrm{Pic}}_{X_m/\underline{C}} \to \underline{\mathrm{Pic}}_{X_n/\underline{C}}
\]
are just vector groups.
\note{dans « $m \geqslant n$ », le « $\geqslant$ » est ajouté à l'encre.}

\page{227}
Parts (i) and (ii) yield the

\emph{Corollary 1} The following conditions are equivalent:

(i) There exists an open $U$ \supplied{$\ni$} $y$ such that $X$
\supplied{$|$} $U$ is projective over $U$

(ii) For every $n$, $X_n$ is a projective analytic space.

For instance, if $\dim X_0 \leqslant 1$, then (ii) and hence (i) holds.

The proof of the theorem only uses Grauerts analogues of the algebraic
theorems of finiteness and comparison for direct images (of his blue paper) and
the usual exact sequence
$0 \to \underline{Z} \to \underline{O} \to \underline{O}^{\bullet} \to 0$,
together with some standard use of Mittag-Leffler story and five lemma. It is
valid in fact for any $H^i(\underline{O}^{\bullet})$, not only $i = 1$ (which
seems the only one however to have geometric significance).
\note{les flèches de la suite exacte manquent à la frappe et sont restituées.}

In the case of a \emph{formal} scheme proper over a complete noeth.\ local
ring with residue field of \emph{caracteristic} $0$, the analogon of the
previous theorem (reducing to statements (i), (i bis)) \struck{\ill{}} hold
true, and I wrote a purely algebraic proof of this, relying only on the fact
that the Kernel and Cokernel of
$\underline{\mathrm{Pic}}_{X_{n+1}}$ \supplied{$\to$}
$\underline{\mathrm{Pic}}_{X_n}$ are without torsion, and Néron's finiteness
theorem; in particular, the analogon of corollary 1 holds true in this case.
These results break down of course in car.~$> 0$.
\note{« formal » est souligné à l'encre ; « $>$ » est ajouté à l'encre.}

However, using the (as yet unwritten !) Gaga of
Serre-Grauert-Remmert-Grothendieck (cf Grauert-Remmerts paper, complemented
by the method of an old talk of mine in Cartan's Seminar, to recover the case
of proper morphisms of schemes from the projective one, via Chow's lemma …),
the analytic theorem above yiels an interesting intrinsic property of analytic
algebras over $\underline{C}$, with respect to algebraic geometry

\page{228}
over such a local ring:

\emph{Theorem} Let $A$ be an analytic algebra over $\underline{C}$ (we can
suppose $A$ to be the ring of convergent power series in $n$ variables),
$\hat{A}$ its completion, $Y$ and $\hat{Y}$ the spectra, $X$ a proper scheme
over $Y$, $\hat{X} = X \times_Y \hat{Y}$. Then:

(i) The inverse system $(\operatorname{Pic}(X_n))$ satisfies MLAR (as stated
above, this depends only on the car.~$0$ assumption for the residue field).

(ii) $\operatorname{Pic}(X)$ and $\operatorname{Pic}(\hat{X})$ have same image
in $\operatorname{Pic}(X_n)$, -- namely the group of “universal images”.
(NB recall $\operatorname{Pic}(\hat{X}) \simeq \varprojlim \operatorname{Pic}(X_n)$ ).

(iii) In order for $\operatorname{Pic}(X) \to \operatorname{Pic}(\hat{X})$ to be
an isomorphism, it is necessary and sufficient that
$\operatorname{supp} R^1 f_*(\underline{O}_X^{*}) \subset (y)$, i.e.\
$H^1(X, \underline{O}_X)$ of finite length over $A$. (NB It amounts also to
the same to ask that the inverse system of the subgroups
$\operatorname{Pic}'(X_n)$ of universal images is constant for large $n$).

We get for example:

\emph{Corollary 1} The following conditions are equivalent: (i) $X/Y$
projective (ii) $\hat{X}/\hat{Y}$ projective (iii) For every $n$, $X_n/Y_n$
projective.

This applies for instance if $\dim X_0 \leqslant 1$.

Now, applying your theorem of \emph{resolution of singularities}, and
Mumford's method of relating the local Picard group of $A$ to the global
Picard group of a regular scheme dominating $A$ birationnally, one gets from
the last statement (iii) of last theorem:

\emph{Corollary} 2. Let $A$ be as above, assume \struck{\ill{}}
$Y' = Y - (y)$ regular, and let $\hat{Y}' = \hat{Y} - (\hat{y})$. Then
$\operatorname{Pic}(Y') \to \operatorname{Pic}(\hat{Y}')$ is an isomorphism.

This explains “à priori” (when $A$ is normal) why Mumford was

\page{229}
\add{able} to introduce \struck{\ill{}} on the group of divisor-classes of $A$
a structure of an analytic group (which in fact is algebraic…), which from
the algebraic point of view should be possible rather for the group of divisor
classes of the completion $\hat{A}$; of course Mumford uses directly the same
kind of argument I used.

I do not know if in the last statement, the hypothesis that $Y'$ is regular
(“$y$ isolated singularity”) is essential; we could dispense with it
\struck{\ill{}} and replace it simply by “$A$ reduced” if you can prove by
your theory of resolution the following: if $f\colon X \to Y$ is proper
“birational”, $X$ regular, then $R^1 f_*(\underline{O}_X) = 0$. I understand
you proved this if $Y$ also is regular (which is easily checked by your
theory), but I wonder if this is \emph{really} needed.
\marginal{of course it is}
I would not be surprised either if in this statement, $Y'$ can be replaced by
any open subset of $Y$ (replacing of course $\hat{Y}'$ by the inverse image of
the latter). Moreover, I would expect the analoguous statements to hold for
$\pi_1$, more generally for all “topological” invariants as Weil homology,
homotopy groups etc, that can be defined for schemes. This should be related
to the fact that all these invariants vanish for the geometric fibers of the
morphism $\operatorname{Spec}(\hat{A})$ \supplied{$\to$}
$\operatorname{Spec}(A)$. This is easy to check at least for $\pi_0$ (and is
true in fact for any henselian ring which is a “good” ring); however I do not
know if this is true also for $\pi_1$.
\note{la note marginale « of course it is », au crayon, est en regard de « really needed », souligné à la main.}

Besides, I would not be surprised if most of the previous results (namely
parts (ii) and (iii) of the second theorem, and the two corollaries, as well
as the \struck{previous} conjectures of the previous section) did hold true for
any “good” ring which is henselian, \add{or} at least for the

\page{230}
“henselian closures” of the local rings arising from algebras of finite type
over a field, or over the integers, -- although I do not have any result along
these lines (except \add{those} stemming from my remark on $\pi_0$).
\struck{\ill{}} This can be stated of course directly in terms of conjectures
for the latter local rings without explicit reference to a henselian closure,
for instance corollary 2 would yield the conjectural statement: Let $A$ be a
local \struck{\ill{}} ring of an algebra of finite type over a field, $\hat{A}$
its completion, $Y$ and $\hat{Y}$ the spectra, $Y'$ and $\hat{Y}'$ the
complements of the closed points, then any invertible sheaf on $\hat{Y}'$ can
be defined by an invertible sheaf on some $Y_1'$, where $Y_1$ is local and
$Y_1$ \supplied{$\to$} $Y$ is étale with trivial residue field extension (i.e.\
inducing an isomorphism for the completions in $A$ \supplied{$\to$} $A_1$). I
wonder what information is given by Mumfords example in his blue paper, p.~16,
which I believe yields a case where the invertible sheaf considered does not
come from an invertible sheaf on $Y'$ ? I was not able to understand his
construction.
\note{« $\pi_0$ » est écrit à l'encre ; dans « on $\hat{Y}'$ », le « $\hat{Y}'$ » est corrigé à la machine sur un « $\hat{X}$ ».}

Anyhow, one should be able to determine whether or not \add{the} analytic
algebras over $\underline{C}$ have any significant intrinsic property which is
not shared by all “good” henselian rings with residu field of car.~$0$
(I recall that by good I mean -- “quotient of a regular local ring $B$ such
that the fibers of $\operatorname{Spec}\hat{B} \to \operatorname{Spec}(B)$ are
universally regular).

Please give my regards to Waka, and also Mireille's ; she just got the parcel
from Waka, and was extremely pleased, in fact, she slipped into her new
bed-shirt on the spot, and is delighted by it in every respect.

Sincerely yours

\subsection{Pic d'un schéma formel en car.~0}
\note{titre encadré en haut à droite de la page 232.}

\page{232}
\textbf{Lemme} Soit $\varphi\colon G \to H$ un homomorphisme de type fini de
groupes loc.\ de type fini sur un corps $k$. Alors

(i) $\underline{\Pi}_0(\varphi)\colon G/G^0 \to H/H^0$ est de type fini, en
particulier son noyau est fini.

\struck{[Si $N = \operatorname{Ker}\varphi$ est connexe, alors
$\underline{\Pi}_0(\varphi)$ est \ill{} \ill{} \ill{}]}
\marginal{(\ill{} si $H/H^0$ \ill{} \ill{} \ill{} \ill{} $G/G^0$)}

(ii) Supposons que $N = \operatorname{Ker}\varphi$ soit connexe, soit
$H' = \operatorname{Im}\varphi$, \add{et $i\colon H' \to H$ l'immersion
canonique.} Alors
\[
\operatorname{Ker}\underline{\Pi}_0(\varphi)
\text{\struck{$= \operatorname{Ker}\underline{\Pi}_0(i)$}}
\simeq \bigl(H'/H'^0\bigr) \cap \bigl(H/H^0\bigr)^0 .
\]
\add{et $\operatorname{Coker}\underline{\Pi}_0(\varphi)$, $\operatorname{Coker}\varphi$}
\add{ne changent} pas quand on remplace $(G, H, \varphi)$ par $(H', H, i)$
ou même par $(H'/H^0, H/H^0, i')$ …
\note{la formule pour $\operatorname{Ker}\underline{\Pi}_0(\varphi)$ est barrée d'un trait qui la traverse en entier ; les ajouts interlinéaires la remplacent par l'énoncé d'invariance. Sous « $H'/H^0$ », une annotation : « discret $\subset F$ ».}

\struck{Corollaire} \add{(iii)} Supposons \add{que $\operatorname{Ker}\varphi$
soit connexe et} que $\operatorname{Coker}\varphi$ soit “géométriquement sans
torsion”. Alors $\underline{\Pi}_0(\varphi)$ est injectif et
$\operatorname{Coker}\underline{\Pi}_0(\varphi)$ est sans torsion.
\note{« (iii) » est cerclé et remplace l'étiquette « Corollaire » biffée.}

Dém.\ \add{de (iii)}. Grâce à (ii) précédent, on est ramené au cas où
$\varphi$ est une immersion fermée, et où $G$ est \uncertain{discret}. Mais
alors $\operatorname{Coker}\varphi$ contient $H^0/G \cap H^0$ où
$G \cap H^0 = \operatorname{Ker}\underline{\Pi}_0(\varphi)$. Si ce quotient
est sans torsion, il faut que, ou bien

\page{234}
a) car $k > 0$, \struck{\ill{}} \add{et} $H^0$ \struck{\ill{}}
\uncertain{réduit} : \ill{} \ill{}
alors $G \cap H^0 = \{e\}$ i.e.\ $\underline{\Pi}_0(\varphi)$ est injectif,
et du \uncertain{point} de vue géométrique \uncertain{on a}
$\varphi = \underline{\Pi}_0(\varphi)$ donc
$\operatorname{Coker}\underline{\Pi}_0(\varphi)$ est sans torsion

b) car $k = 0$, $H^0$ est un groupe \emph{vectoriel}, \struck{et} donc le
sous-groupe fini $G \cap H^0 = 0$. Donc \struck{\ill{}} on a une suite exacte
\[
0 \to H^0 \to H/G \to \operatorname{Coker}\underline{\Pi}_0(\varphi) \to 0 .
\]
qui splitte du point de vue géométrique [car $H^0$ est divisible]. Comme $H/G$
est sans torsion, il en est donc de même de
$\operatorname{Coker}\underline{\Pi}_0(\varphi)$.
\marginal{$G \subset H$}
\note{sous « $\operatorname{Coker}\underline{\Pi}_0(\varphi)$ », écrit au-dessus d'un mot barré, l'égalité « $= H/(G \cdot H^0)$ ».}

\emph{Corollaire 1} Soit $(G_n)_{n \in \mathbf{Z}}$ un système projectif de
groupes \add{commutatifs} \struck{\ill{}} loc.\ de type fini sur le corps $k$.
Supposons \add{a)} \struck{\ill{}} $G_0/G_0^0$ est à engendrement fini
\add{b)} pour tout entier $n$, l'homomorphisme de transition
$G_{n+1} \to G_n$ \add{est de type fini,} \struck{a} a un noyau connexe, et un
conoyau qui est géométriquement \struck{sans} sans torsion

\add{b)} \uncertain{pour tout} $n$, $\underline{\Pi}_0(G_n)$ est à engendrement fini.
\note{la dernière ligne, serrée en bas de page sous la précédente, porte elle aussi l'étiquette « b) » ; on ne voit pas si elle remplace la condition b) ou si elle commence déjà la conclusion, qui se poursuit page 236.}

\page{236}
Sous ces conditions

(i) Le système projectif $(\underline{\Pi}_0(G_n))_n$ est essentiellement
constant.

(ii) Les systèmes $(G_n^0)_n$ [donc aussi le système $(G_n)$] satisfont
(M.L.).

(iii) \struck{\ill{}} \add{(\ill{} \ill{} \ill{})} Il existe un $n_0$ tel que
pour $n \geqslant n_0$, le \struck{noyau et conoyau} noyau de
$G_n \to G_{n_0}$ soit un groupe vectoriel en car.~$= 0$, et
\struck{sans torsion [dans} \struck{presque} radiciel en car.~$> 0$.
\struck{[Si la car.\ est nulle et si \ill{} les $(G_{n+1} \to G_n)$ \ill{}
\ill{}, alors les noyaux des $(G_{n+1} \to G_n)$ sont des groupes
vectoriels,}
\marginal{(\ill{} pour $n \geqslant n_0$ \ill{})}

\emph{Corollaire 2} Supposons que les
$\operatorname{Ker}(G_{m+1} \to G_m)$ soient \struck{\ill{}} \add{\ill{}}
\add{unipotents} \struck{\ill{}} groupes \struck{unipotents connexes},
\struck{supposons $k$ parfait et} supposons $k$ parfait
[le cas important étant le car.~$0$]. Alors le système des
$(G_n(k))_n$ satisfait ML.

\page{238}
\note{en tête de page, un brouillon : « $G_m$ \ill{} $P \to Q \to R$ », avec \struck{$\operatorname{Im} P \cap \operatorname{Ker} Q$} remplacé par $\operatorname{Im}(P \to Q) \cap \operatorname{Ker}(Q \to R)$.}

\struck{Dém.} Si \ill{}, si $G_{n,m} = \operatorname{Im}(G_m \to G_n)$ pour
$m \geqslant n$, on doit prouver que
\[
G_{n,m}(k) = \operatorname{Im}(G_m(k) \to G_n(k))
\]
\uncertain{Il suffit} \uncertain{de voir} que $H^1(k, N_{m,n}) = 0$. Par
suite, \add{\ill{} \ill{},} le \uncertain{système} des
$\operatorname{Im}(G_m(k) \to G_n(k))$ \uncertain{sont} stationnaires.

\emph{N.B.} Si les $G_n$, \ill{} pour \uncertain{alg.}, \ill{} \ill{}
\add{\ill{} unipotent \uncertain{connexe}}, \uncertain{donc} si les
$\operatorname{Ker}(G_{n+1} \to G_n)$ sont des groupes
\uncertain{unipotents}, \struck{\ill{} \ill{}} les
$\operatorname{Ker}(G_m \to G_n)$ \uncertain{pour} \ill{} sont des groupes
vectoriels, en \uncertain{particulier} \ill{} \ill{}.
\marginal{\ill{}}
\note{le N.B.\ est tenu par un trait vertical dans la marge gauche ; plusieurs mots y sont surchargés.}

\textbf{Th.} Soit $\mathfrak{X}$ un schéma formel \add{propre} sur un
\ill{} local noethérien complet $A$, \uncertain{contenant} \uncertain{un
corps} de car.~$0$, tel que $k(A)$ soit fini sur $k$. Considérons les $X_n$
comme des schémas propres sur $k$, considérons le syst.\ projectif
\[
\bigl(\underline{\mathrm{Pic}}_{X_n/k}\bigr) .
\]
Alors ce syst.\ satisfait aux conditions

\page{240}
du Cor.~1 précédent, et le système des
$\underline{\mathrm{Pic}}_{X_n/k}(k)$ satisfait à M.L.

\marginal{\emph{N.B.} On \uncertain{oublie} \ill{} $\Rightarrow$ M.L.A.R.}

\emph{N.B.} On suppose l'existence des $\underline{\mathrm{Pic}}_{X_n/k}$,
\ill{} prouvée en général \uncertain{par} \emph{Murre}. Dans le cas
projectif, \uncertain{elle} \ill{} \ill{} par Grothendieck -- Oort.
\uncertain{Le fait} que les
$\underline{\mathrm{Pic}}_{X_{n+1}/k} \to \underline{\mathrm{Pic}}_{X_n/k}$
sont de type fini, et \struck{\ill{}} le noyau \struck{\ill{}}
\uncertain{connexe} \add{(\uncertain{dont} \ill{} becomes \ill{})}
\struck{\ill{}}, le conoyau sans \uncertain{torsion}
\add{(becomes \ill{})}, \uncertain{se voit} sur la théorie de Oort
\struck{des obstructions} [cf \ill{} ex.\ \ill{} de \uncertain{cohomologie}].
De plus, on utilise le \add{\uncertain{théorème} \ill{}} \ill{} de
\emph{Néron} \struck{\ill{}} pour $\underline{\mathrm{Pic}}_{X_0/k}$.
\note{passage très surchargé : deux insertions entre parenthèses, en partie en anglais (« becomes »), sont reliées par des accolades ; leur place exacte dans la phrase est incertaine.}

(Soit $\mathfrak{X}$ \uncertain{propre} \ill{} \uncertain{loc.\ noeth.\
complet} \ill{} \ill{} \uncertain{résiduel} $k$ de car.~$0$ \ill{} \ill{}
\ill{} \ill{} \ill{} \ill{} \ill{} \ill{}
$(\underline{\mathrm{Pic}}(X_n))_{n \in \mathbf{Z}}$. Alors \ill{}.)
\note{note entre parenthèses, écrite serrée sur quatre lignes au-dessus du corollaire, largement illisible.}

\emph{Corollaire} \ill{} à M.L.

\ill{} \ill{} \uncertain{Lorsque} \ill{} \ill{} \ill{}
$k \subset A$, $X_0$ \ill{} \ill{} \ill{} \ill{} $/k$, \ill{} \ill{} \ill{}
\note{les dernières lignes sont barrées de deux traits obliques et ne se lisent qu'en fragments ; l'argument se poursuit au-delà du lot.}

\end{document}
